\chapter{SageMath Basics for Post-Quantum Cryptography}
\label{ch:sagemath}

\section{Why this chapter matters}
This chapter is not meant to teach SageMath in isolation. Its purpose is to build a small experimental laboratory for post-quantum cryptography. Later chapters on LWE, RLWE, Kyber, Dilithium, McEliece, and lattice attacks all depend on the objects introduced here.

\section{Learning goals}
After finishing this chapter, you should be able to:

\begin{itemize}
  \item create integers, rationals, reals, complex numbers, finite fields, vectors, matrices, and polynomials in SageMath;
  \item distinguish Sage objects from ordinary Python values;
  \item express algebraic structures such as $\mathbb{F}_q$, $R_q = \mathbb{Z}_q[x]/(f(x))$, and $Ax=b$ in Sage;
  \item run small experiments that foreshadow later cryptographic constructions.
\end{itemize}

Every algebraic object used later in the book is a specialization of a small handful of number systems. Figure~\ref{fig:number-systems} lays out that map before we start building each piece in Sage.

\begin{figure}[H]
\centering
\begin{tikzpicture}[node distance=9mm and 9mm]
  \node[flownode] (zz) {$\mathbb{Z}$\\[-2pt]\scriptsize \texttt{ZZ}};
  \node[flownode, right=of zz] (qq) {$\mathbb{Q}$\\[-2pt]\scriptsize \texttt{QQ}};
  \node[flownode, right=of qq] (rc) {$\mathbb{R}/\mathbb{C}$\\[-2pt]\scriptsize \texttt{RR/CC}};
  \node[flownode, below=of zz] (gf) {$\mathrm{GF}(q)$\\[-2pt]\scriptsize finite field};
  \node[flownode, below left=15mm and 6mm of gf] (vecs) {Vectors/Matrices\\[-2pt]\scriptsize over $\mathrm{GF}(q)$};
  \node[flownode, below right=15mm and 6mm of gf] (rq) {$R_q=\mathrm{GF}(q)[x]/(f(x))$};
  \draw[flowarrow] (zz) -- (qq);
  \draw[flowarrow] (zz) -- (rc);
  \draw[flowarrow] (zz) -- (gf);
  \draw[flowarrow] (gf) -- (vecs);
  \draw[flowarrow] (gf) -- (rq);
\end{tikzpicture}
\caption{The map of algebraic objects this book builds on. Ordinary integers $\mathbb{Z}$ extend to fractions $\mathbb{Q}$ and to $\mathbb{R}/\mathbb{C}$ for numerical and Fourier-style reasoning, but the object cryptography actually computes with is the finite field $\mathrm{GF}(q)$, obtained by reducing $\mathbb{Z}$ modulo a prime $q$. Everything in later chapters -- vectors and matrices for LWE, and the quotient ring $R_q$ for RLWE and Kyber -- is built on top of $\mathrm{GF}(q)$.}
\label{fig:number-systems}
\end{figure}

\section{Installing SageMath}

\subsection{Option 0: install nothing}
Every code listing in this book runs in a browser at
\href{\playgroundurl}{appliedpqc.io/playground.html}, on the free SageMath Cell
service. Nothing is installed, and each listing arrives with the earlier
listings of its chapter already run, so any snippet can be tried on its own.
The one thing a browser cannot do is the slowest work --- Falcon key
generation, and signing under the SLH-DSA \texttt{s} parameter sets --- which
needs a local Sage.

\subsection{Option 1: conda-forge}
The quickest local install on any platform, and the one that gives the newest
release:

\begin{lstlisting}[language=bash]
conda create -n sage -c conda-forge sage
conda activate sage
\end{lstlisting}

The package is named \texttt{sage}. Note there is no \texttt{sagemath-standard}
package on conda-forge, despite the name appearing in some older instructions.

\subsection{Option 2: Docker}
To keep the host clean, start SageMath from a container:

\begin{lstlisting}[language=bash]
docker run -it sagemath/sagemath sage
\end{lstlisting}

Once the shell starts, you can enter Sage commands directly. To run a script
rather than work interactively, override the image's entry point, and mount a
\emph{writable} directory --- Sage writes a preparsed \texttt{.sage.py} beside
each source file:

\begin{lstlisting}[language=bash]
docker run --rm --entrypoint bash -v "$PWD:/home/sage/work" \
    -w /home/sage/work sagemath/sagemath:10.9 -lc 'sage myscript.sage'
\end{lstlisting}

\subsection{Option 3: system packages}
On Debian and Ubuntu:

\begin{lstlisting}[language=bash]
sudo apt install sagemath
\end{lstlisting}

On macOS, SageMath is distributed as a cask rather than a formula:

\begin{lstlisting}[language=bash]
brew install --cask sage
\end{lstlisting}

On Windows, install WSL2 and then use conda-forge or \texttt{apt} inside the
Linux environment. Debian and Ubuntu currently ship Sage 9.5, some way behind
conda-forge; everything in this book works on 9.5 or newer.

\subsection{Notebook workflow}
For experiments and exploratory work, a notebook interface is often most convenient:

\begin{lstlisting}[language=bash]
sage -n jupyter
\end{lstlisting}

Then open the displayed local URL in a browser.

\section{SageMath and Python}
SageMath is built on top of Python, but it adds richer mathematical types and objects. A plain Python integer and a Sage integer are not the same thing.

\begin{lstlisting}[language=Python]
a = 3
b = 4
a + b
\end{lstlisting}

In Python, the result is an ordinary integer. In Sage, the corresponding object may carry algebraic meaning.

\begin{lstlisting}[language=Python]
a = Integer(3)
b = Integer(4)
a + b
\end{lstlisting}

To inspect the type, use:

\begin{lstlisting}[language=Python]
type(a)
\end{lstlisting}

This matters in cryptography because the symbol $3$ may represent an ordinary integer in one context, but a field element in another. For example, $3 \in \mathrm{GF}(7)$ is not the same object as the integer $3$ in $\mathbb{Z}$.

\section{Basic number systems}
SageMath provides several standard number domains that are useful in later cryptographic work.

\subsection{Integers $\mathbb{Z}$}
The integer ring is available as \texttt{ZZ}.

\begin{lstlisting}[language=Python]
ZZ
\end{lstlisting}

You can create integers explicitly:

\begin{lstlisting}[language=Python]
a = ZZ(10)
a
\end{lstlisting}

Basic arithmetic works as expected:

\begin{lstlisting}[language=Python]
a = ZZ(10)
b = ZZ(3)
a + b
\end{lstlisting}

Division behaves more carefully than in Python because Sage keeps symbolic and exact arithmetic whenever possible:

\begin{lstlisting}[language=Python]
10/3
\end{lstlisting}

The result is the rational number $\frac{10}{3}$. Integer division is available through:

\begin{lstlisting}[language=Python]
10//3
\end{lstlisting}

and modular arithmetic through:

\begin{lstlisting}[language=Python]
10 % 3
\end{lstlisting}

\subsection{Rationals $\mathbb{Q}$}
The rational field is \texttt{QQ}.

\begin{lstlisting}[language=Python]
a = QQ(1)/3
a
\end{lstlisting}

This is useful because many lattice algorithms, such as LLL and Gram--Schmidt-style computations, work naturally over rational numbers.

\subsection{Reals $\mathbb{R}$ and complexes $\mathbb{C}$}
SageMath also supports real and complex domains.

\begin{lstlisting}[language=Python]
RR
CC
\end{lstlisting}

For example:

\begin{lstlisting}[language=Python]
sqrt(2)
RR(sqrt(2))
CC(1 + 2*I)
\end{lstlisting}

These domains become relevant later when discussing numerical methods, Fourier-style transforms, and related implementation ideas.

\section{Finite fields}
Finite fields are among the most important structures in modern cryptography. A finite field is usually written as $\mathrm{GF}(q)$ or $\mathbb{F}_q$. Before computing with them, we define them precisely; they underlie every construction in this book.

\begin{definition}[Field]
A \emph{field} is a set $\mathbb{F}$ with addition and multiplication such that $(\mathbb{F},+)$ and $(\mathbb{F}\setminus\{0\},\cdot)$ are both abelian groups (with identities $0$ and $1$) and multiplication distributes over addition. Concretely: one can add, subtract, multiply, and divide by anything nonzero.
\end{definition}

A \emph{finite} field is simply a field with finitely many elements. Their existence is completely classified.

\begin{theorem}[Classification of finite fields]
A finite field of order $q$ exists if and only if $q=p^m$ is a power of a prime $p$, and it is then unique up to isomorphism, denoted $\mathbb{F}_q$ or $\mathrm{GF}(q)$. When $m=1$ it is simply the integers modulo $p$,
\[
\mathbb{F}_p = \mathbb{Z}/p\mathbb{Z} = \{0,1,\dots,p-1\},
\]
with arithmetic reduced modulo $p$.
\end{theorem}

Almost every modulus in this book is prime, so the working picture is the clean one: $\mathbb{F}_q$ is ``clock arithmetic'' modulo $q$ in which division is also available. This field structure makes tools such as Gaussian elimination available when they are needed. It is not, however, why ML-KEM or ML-DSA key generation works: those algorithms sample secrets and public values rather than solving a random linear system. Their prime moduli are chosen for several interacting reasons, including algebraic structure, NTT compatibility, noise distributions, correctness, security, and implementation efficiency.

\subsection{Creating a finite field}
For example, the field with seven elements is:

\begin{lstlisting}[language=Python]
F = GF(7)
F
\end{lstlisting}

Elements are created by coercing values into the field:

\begin{lstlisting}[language=Python]
a = F(3)
a
\end{lstlisting}

Arithmetic is performed modulo the prime:

\begin{lstlisting}[language=Python]
a = F(6)
a + F(3)
\end{lstlisting}

The result is $2$, because $6 + 3 = 9 \equiv 2 \pmod{7}$.

\subsection{Inverses}
In cryptography, inverses appear constantly. For a field element $a$, the inverse $a^{-1}$ satisfies $a a^{-1} = 1$.

\begin{lstlisting}[language=Python]
a = F(3)
a.inverse()
\end{lstlisting}

Since $3 \cdot 5 = 15 \equiv 1 \pmod{7}$, the inverse is $5$.

Under the hood, \texttt{a.inverse()} does not search through all field elements; it runs the extended Euclidean algorithm, which computes $\gcd(a,q)$ while simultaneously tracking the coefficients of a Bézout identity $as+tq=\gcd(a,q)$, so that when $\gcd(a,q)=1$ the coefficient $s$ is exactly $a^{-1} \bmod q$.

\begin{algorithm}[H]
\caption{Extended Euclidean Algorithm for Modular Inverse}
\label{alg:ext-euclid}
\begin{algorithmic}[1]
\Require Integer $a$ with $1 \le a < q$; modulus $q$ with $\gcd(a,q)=1$
\Ensure $a^{-1} \bmod q$, the unique $x \in \{0,\dots,q-1\}$ with $ax \equiv 1 \pmod{q}$
\State $(r_0, r_1) \gets (q, a)$
\State $(s_0, s_1) \gets (0, 1)$ \Comment{invariant: $a\, s_i \equiv r_i \pmod q$ at every step}
\While{$r_1 \neq 0$}
    \State $t \gets \lfloor r_0 / r_1 \rfloor$
    \State $(r_0, r_1) \gets (r_1,\ r_0 - t\, r_1)$
    \State $(s_0, s_1) \gets (s_1,\ s_0 - t\, s_1)$
\EndWhile
\State \Return $s_0 \bmod q$ \Comment{$r_0=\gcd(a,q)=1$, so $a\,s_0 \equiv 1 \pmod q$}
\end{algorithmic}
\end{algorithm}

\section{Linear algebra over a finite field}
Finite-field linear algebra is the starting point for many lattice-based constructions.

\begin{lstlisting}[language=Python]
F = GF(7)
A = Matrix(F, [[1, 2], [3, 4]])
A
\end{lstlisting}

The inverse can be computed as:

\begin{lstlisting}[language=Python]
A.inverse()
\end{lstlisting}

and one can verify the identity matrix:

\begin{lstlisting}[language=Python]
A * A.inverse()
\end{lstlisting}

This is the algebraic setting behind equations such as $A s + e = b$, which will appear later in the study of LWE.

\section{Vectors and matrices}
Vectors are convenient for representing secret values and noise terms.

\begin{lstlisting}[language=Python]
v = vector(F, [1, 2])
v
\end{lstlisting}

Matrix-vector multiplication is then straightforward:

\begin{lstlisting}[language=Python]
A * v
\end{lstlisting}

The output is \texttt{(5, 4)}. The vector has length $2$, matching the two columns of $A$; dimensions must agree for matrix--vector multiplication.

\section{Polynomial rings}
Polynomial rings are essential for the ring-based versions of lattice problems, especially RLWE and related constructions.

\subsection{Building a polynomial ring}

\begin{lstlisting}[language=Python]
F = GF(7)
R.<x> = PolynomialRing(F)
\end{lstlisting}

Now $x$ is a formal variable in the ring.

\begin{lstlisting}[language=Python]
f = x^3 + 2*x + 1
f
\end{lstlisting}

Addition and multiplication are performed in the usual algebraic way:

\begin{lstlisting}[language=Python]
f + f
f * f
\end{lstlisting}

\subsection{Quotient rings}
A quotient ring is created by imposing a polynomial relation. For example:

\begin{lstlisting}[language=Python]
S.<x> = QuotientRing(R, R.ideal(x^3 + 1))
\end{lstlisting}

In this quotient, the relation $x^3 = -1$ holds. This idea is closely related to the ring structure used in Kyber, where one often works in a ring of the form

\[
R_q = \mathbb{Z}_q[x]/(x^{256}+1).
\]

\section{A first cryptographic experiment}
As a first experiment, we can verify the behavior of modular inverses.

\begin{lstlisting}[language=Python]
F = GF(17)
for i in range(1, 17):
    a = F(i)
    print(i, a.inverse() * a)
\end{lstlisting}

Every value should evaluate to $1$ in the field.

\section{A first LWE-style example}
To foreshadow later chapters, we can create a simple linear system over a finite field:

\begin{lstlisting}[language=Python]
q = 17
F = GF(q)
A = random_matrix(F, 3, 3)
s = random_vector(F, 3)
b = A * s
A
s
b
\end{lstlisting}

This is the basic algebraic shape behind LWE before noise is introduced.

\section{Exercises}
Try to solve the following exercises on your own.

\begin{enumerate}
  \item Create $\mathrm{GF}(13)$ and compute the inverse of every nonzero element.
  \item Create $R = \mathrm{GF}(17)[x]$ and compute $(x^2 + 3x + 1)^2$.
  \item Create $\mathrm{GF}(17)[x]/(x^4 + 1)$ and verify that $x^4 = -1$.
  \item Repeatedly sample $A \in \mathrm{GF}(17)^{4 \times 4}$ until \texttt{A.is\_invertible()} is true, then verify that $A^{-1}A = I$. What happens for a singular sample?
\end{enumerate}

\section{Connection to the next chapters}
The next chapter, \emph{Linear Algebra Foundations}, will deepen the linear-algebra viewpoint needed for lattice problems. After that, the polynomial viewpoint will be developed more fully in \emph{Polynomial Rings}. These two chapters prepare the ground for LWE, RLWE, and ultimately Kyber and Dilithium.
